% \documentclass[letterpaper,journal]{IEEEtran}
\documentclass[letterpaper,journal]{IEEEtran}
\usepackage{amsmath,amsfonts}
\usepackage{algorithmic}
\usepackage{algorithm}
\usepackage{array}
\usepackage[caption=false,font=normalsize,labelfont=sf,textfont=sf]{subfig}
\usepackage{textcomp}
\usepackage{stfloats}
\usepackage{url}
\usepackage{verbatim}
\usepackage{graphicx}
\usepackage{cite}
\usepackage{mathptmx} % use matching math font
\usepackage{newtxtext} % use Times font with small caps and italic support
\usepackage{xcolor}
\usepackage{comment}
\usepackage{amsmath}
\usepackage{amsfonts}
\usepackage{enumitem}
\usepackage{underscore}
\usepackage{orcidlink}
\newcommand{\authororcid}[2]{#1\,\orcidlink{#2}}
\usepackage{tabu}                   
\usepackage{booktabs}           

\hyphenation{op-tical net-works semi-conduc-tor IEEE-Xplore}
% updated with editorial comments 8/9/2021

\RequirePackage[nameinlink,capitalise]{cleveref}
\crefname{table}{Tab.}{Tabs.}
\Crefname{table}{Table}{Tables}
\crefname{section}{Sec.}{Secs.}
\Crefname{section}{Section}{Sections}
\crefname{subsection}{Sec.}{Secs.}
\Crefname{subsection}{Section}{Sections}
\crefname{subsubsection}{Sec.}{Secs.}
\Crefname{subsubsection}{Section}{Sections}
\crefname{appendix}{Appx.}{Appx.}
\Crefname{appendix}{Apendix}{Appendices}
\crefname{subappendix}{Appx.}{Appx.}
\Crefname{subappendix}{Apendix}{Appendices}
\crefname{subsubappendix}{Appx.}{Appx.}
\Crefname{subsubappendix}{Apendix}{Appendices}

\begin{document}

\title{From Interaction to Systems Thinking: A Visual Analytics Approach for Exploring Pupils' Behavior in a Digital Learning Environment}

\author{\authororcid{Peilin Yu}{0000-0003-4820-6123},
        \authororcid{Aida Nordman}{0000-0002-9601-5981},
        \authororcid{Mina Mani}{0000-0003-2669-8059},
        \authororcid{Konrad Sch\"{o}nborn}{0000-0001-8888-6843},\\
        \authororcid{Marta Koc-Januchta}{0000-0001-6313-475X},
        \authororcid{Gunnar H\"{o}st}{0000-0003-1032-2145},
        \authororcid{M{\aa}ns Gezelius}{0009-0005-4038-842X},
        \authororcid{Lonni Besan{\c c}on}{0000-0002-7207-1276},
        and \authororcid{Katerina Vrotsou}{0000-0003-4761-8601}
% \author{Peilin Yu,
%         Aida Nordman,
%         Mina Mani,
%         Konrad Sch\"{o}nborn,\\
%         Marta Koc-Januchta,
%         Gunnar H\"{o}st,
%         M{\aa}ns Gezelius,
%         Lonni Besan{\c c}on,
%         and Katerina Vrotsou
        % <-this % stops a space
\thanks{Peilin Yu, Aida Nordman, Mina Mani, Konrad Sch\"{o}nborn, Marta Koc-Januchta, Gunnar H\"{o}st, M{\aa}ns Gezelius, Lonni Besan{\c c}on, and Katerina Vrotsou are with Link\"{o}ping University. E-mail: \{first\_name.last\_name\}@liu.se.}% <-this % stops a space
% \thanks{Manuscript received April 19, 2021; revised August 16, 2021.}
}

% The paper headers
\markboth{}%
{Yu \MakeLowercase{\textit{et al.}}: From Interaction to Systems Thinking: A Visual Analytics Approach for Exploring Pupils' Behavior in a Digital Learning Environment}

% \IEEEpubid{0000--0000/00\$00.00~\copyright~2021 IEEE}
% Remember, if you use this, you must call \IEEEpubidadjcol in the second column for its text to clear the IEEEpubid mark.

\maketitle


\begin{figure*}[!t]
    \centering
    \includegraphics[width=\textwidth]{figs/teaser.pdf}
  \caption{
  Overview of the two main components of \textsc{BEHAVE}.
  The \textsc{Motif and Comparison Panel} features 
  (1) the \textit{Motif Library}, containing all discovered motifs, and 
  (2) a compact visualization of mutually exclusive ST-related events for all interaction sequences, displayed in \textit{durational mode} for cross-sequence comparison. 
  The \textsc{Individual Panel} contains details for a selected pupil.
  It features 
  (3) multidimensional time series summarizing behavioral tendencies;
  (4) a Marey-chart-based representation of interaction events;
  (5) a representation of behavioral building blocks; 
  and
  (6) a compact ST-related event sequence visualization in \textit{sequential mode}.
  The first sequence in the \textsc{Comparison Panel} corresponds to the sequence displayed in the \textsc{Individual Panel}, as indicated by the annotated arrow on the right.
    }
    \label{fig:teaser}
\end{figure*}


\begin{abstract}
We introduce a visual analytics (VA) approach grounded in a systems thinking (ST) framework to explore pupils' problem-solving behaviors and ST reasoning from interaction log data collected in a digital learning environment.
Understanding pupils' interaction behaviors in such environments is essential for uncovering exploration patterns and task-solving strategies, and for informing the design of interactive learning tools.
However, such interaction logs are often large, heterogeneous, and variable in length, making them difficult to analyze, particularly when the goal is to relate observable behaviors to reasoning processes such as ST.
To address this challenge, we introduce an approach that transforms interaction sequences into multidimensional time series that capture temporal progression across behavioral dimensions, and then applies motif discovery, similarity analysis, and clustering to identify recurring behaviors and strategy-aligned cohorts, supported by multi-level visual representations for cohort-, pattern-, and individual-level analysis.
We implement this approach in \textsc{BEHAVE} (\textbf{B}ehavioral \textbf{E}xploration and \textbf{H}ierarchical \textbf{A}nalytics in \textbf{V}isual \textbf{E}nvironments), a system we developed for exploratory analysis.
Our contributions include an ST-grounded VA approach, a motif-based analysis pipeline for interaction logs, the \textsc{BEHAVE} system itself, and a case study demonstrating its usefulness in uncovering behavioral patterns and generating hypotheses about ST reasoning, including Analysis, Synthesis, and Implementation.
We evaluate and validate \textsc{BEHAVE} through a series of workshops with domain experts and a digital learning environment developer using interaction log data from pupils' use of the \textit{Tracing Carbon} interactive learning environment.
Our results show that \textsc{BEHAVE} helps experts identify distinct strategy profiles and uncover behavioral patterns, supporting hypothesis generation about ST reasoning, and informing pedagogical strategies.
\end{abstract}

\begin{IEEEkeywords}
Interaction logs, visual analytics, motif discovery, systems thinking.
\end{IEEEkeywords}

\section{Introduction}

Understanding how pupils engage with interactive digital learning environments is essential to uncover problem-solving strategies and interpret learning processes, as well as to inform the design of more effective learning environments~\cite{wang2023applying}.
Such environments have become increasingly prevalent in science and STEM education, and the analytics tools built around them have evolved in parallel, rapidly progressing from descriptive activity summaries to predictive models, adaptive feedback platforms, and interactive learning dashboards~\cite{siemens2013learning, selwyn2020reimagining, paulsen2024learning, mohseni2024visual}.
As these environments are adopted, they capture interaction data of increasing volume and complexity, creating both an opportunity and a need for methods that can turn rich, fine-grained behavioral traces into pedagogically meaningful insights~\cite{vieira2018visual}.

When analyzed at scale, interaction logs captured in such environments reveal patterns invisible to manual observation, such as systematic problem-solving strategies and recurring behavioral tendencies.
However, interaction logs are often large and heterogeneous, and the gap between low-level logged actions and higher-order cognitive processes remains difficult to bridge~\cite{bogarin2018survey}.
This challenge is particularly acute for open-ended tasks that target reasoning skills, where observable actions obscure the reasoning that produced them.
A quintessential example of such a skill is systems thinking (ST), which is increasingly recognized as a core competency in science education, particularly in topics such as the carbon cycle, where pupils must reason across multiple levels of abstraction simultaneously~\cite{hmelor2007fish}.
ST comprises a range of competencies related to identifying the components of a system, understanding how they are interconnected and change over time, and explaining how higher-order patterns emerge from local interactions~\cite{assaraf2005development, verhoeff2018theoretical}.
Its non-linear, emergent, and feedback-driven reasoning makes it especially difficult to interpret directly from behavioral data.
This makes ST a demanding and informative test case, since any approach relating interaction logs to ST reasoning must confront the log-to-cognition gap.

We therefore introduce a visual analytics (VA) approach that couples analysis of interaction log data with ST. 
The proposed approach is grounded in core ST concepts and reasoning hierarchies, enabling the exploration of pupils' problem-solving strategies, their possible relation to ST reasoning levels, and conceptual difficulties that are difficult to discern from high-level summaries alone.
As a concrete testbed, we use the \textit{Tracing Carbon}~(TC) digital learning environment, designed to help pupils in grades $7$--$9$ learn about the carbon cycle~\cite{mani2026designing}. 
TC suits this study for three reasons: its central task, in which pupils draw arrows to connect carbon reservoirs, produces relational interaction data that maps naturally to ST constructs of connection, cycle, and feedback; it logs interactions at fine granularity, capturing timestamps, source and target reservoirs, and correctness; and it was intentionally designed around the ST hierarchy, aligning the environment with our analytic lens.
Our VA approach transforms raw interaction sequences into multidimensional time series, discovers recurring behavioral motifs between pupils, and organizes the findings through three coordinated panels, each corresponding to a level of the ST hierarchy.
We implement the approach in the prototype system \textsc{BEHAVE} (\textbf{B}ehavioral \textbf{E}xploration and \textbf{H}ierarchical \textbf{A}nalytics in \textbf{V}isual \textbf{E}nvironments) (\cref{fig:teaser}).
Following a design-study methodology~\cite{sedlmair2012design}, \textsc{BEHAVE} was developed through an iterative co-design process with domain experts and evaluated using a case study and structured expert validation based on interaction log data from pupils solving a task in TC.
Because the goal is to support open-ended, exploratory hypothesis generation about pupils' ST reasoning rather than to benchmark task performance, we adopted an insight-oriented qualitative evaluation approach appropriate for this class of tool~\cite{lam2012empirical}.
While we demonstrate the approach on TC, the underlying workflow, which combines theory-informed behavioral abstractions with multi-level sequence analysis, is not specific to this environment and can be transferred to other contexts where interactions can be mapped to a reasoning framework.
The contributions are summarized as follows:

\begin{itemize}[nosep]
    \item \textbf{ST-grounded VA approach}.
    We introduce a VA approach that structures exploration of pupil interaction logs into three ST-related levels: events, recurring patterns, and underlying structures, using ST both as a design lens and an interpretive framework.

    \item \textbf{Tailored motif discovery}.
    We transform discrete interaction sequences into multidimensional time series and identify cross-pupil consensus motifs as candidate recurring behavioral patterns.

    \item \textbf{Interactive system}.
    We propose a VA prototype system that implements the proposed approach, supporting cluster-driven and motif-driven analytical workflows that converge on fine-grained individual inspection and enable hypothesis-driven exploration of the relationship between interactions and ST reasoning.

    \item \textbf{Case study}.
    We demonstrate the usefulness and applicability of the proposed approach using pupil interaction data from the TC environment, revealing previously unseen behavioral patterns and generating new hypotheses about the relationship between problem-solving strategies and ST reasoning.
    The resulting behavioral findings were validated by domain experts through comparison with a previous analysis of the same dataset.
\end{itemize}




%%%%%%%%%%%%%%%%%%%%%%%%%%%%%%%%%%%%%%%%%%%%%%%%%%%%%%%%%%%%%%%%
% DOMAIN
%%%%%%%%%%%%%%%%%%%%%%%%%%%%%%%%%%%%%%%%%%%%%%%%%%%%%%%%%%%%%%%%

\section{Domain Background and Motivation}
\label{sec:domain}
We introduce the digital learning environment that motivates this work and describe ST as a core science competency.

% %%%%%%%%%%%%%%%%%%%%%%%%%%%%%%%%%%%%%%%%%%%%%%%%%%%%%%%%%%%% %
\subsection{Digital Learning Environment: Tracing Carbon}

\textit{Tracing Carbon}~(TC)~\cite{mani2026designing} is a web-based learning environment designed for pupils in grades $7$--$9$ in the Swedish school context (see~\cref{fig:tcInterface}). 
It comprises several interactive tasks intended to support learning about the carbon cycle.
We consider one of the most complex tasks, in which pupils draw arrows representing carbon transfer between carbon reservoirs in the Earth's global carbon cycle.
This task makes TC particularly well-suited to our study for three reasons.
First, it produces \textbf{relational interaction data}: each arrow encodes a directed connection between system components, and sequences of such connections map naturally onto ST constructs of connection, cycle, and feedback, suiting motif-based analysis.
Second, interactions are captured through \textbf{fine-grained logging} of timestamps, mouse actions, source and target reservoirs, and arrow correctness, providing sufficient resolution for sequence- and motif-level analysis.
Third, TC was intentionally \textbf{designed around the ST hierarchy}~\cite{julsgard2026supporting, mani2026designing}, so the environment and our analytic lens are aligned, allowing observed behaviors to be interpreted within an established reasoning framework rather than as isolated actions.
Together, these properties make the task a rich setting for studying how pupils approach the problem, whether they focus on isolated parts or broader relations, and how they organize their reasoning. 

\begin{figure}
    \centering
    \includegraphics[
    width=.93\columnwidth,
    alt={
    A screenshot of the Tracing Carbon digital learning environment. 
    The interface presents twelve interconnected carbon reservoirs across the atmosphere, terrestrial surface, oceans, and fossil fuels. 
    Pupils interact by drawing arrows between reservoirs to represent carbon transfer processes, such as photosynthesis, respiration, and decomposition, with directional links indicating the flow and relative speed (fast versus very slow). 
    The filled yellow circle highlights an active interaction, indicating the pupil's cursor as it creates an arrow link between components. 
    The system logs these interactions to capture pupils' construction of relationships in the carbon cycle and their problem-solving behavior.
    }
    ]{figs/TC_m1c2.pdf}
    \caption{
    A screenshot of the TC digital learning environment. 
    The interface presents $12$ interconnected carbon reservoirs across the atmosphere, terrestrial surface, oceans, and fossil fuels. 
    Pupils interact by drawing arrows between reservoirs to represent carbon transfer processes, such as photosynthesis, respiration, and decomposition, with directional links indicating the flow and relative speed (fast versus very slow). 
    The filled yellow circle highlights an active interaction, indicating the pupil's cursor as it draws a link between components. 
    The system logs these interactions to capture pupils' construction of relationships in the carbon cycle and their problem-solving behavior.
    }
    \label{fig:tcInterface}
\end{figure}



% %%%%%%%%%%%%%%%%%%%%%%%%%%%%%%%%%%%%%%%%%%%%%%%%%%%%%%%%%%%% %
\subsection{Systems Thinking in Science Education}

In STEM education, understanding the carbon cycle is essential for making scientifically informed decisions about climate issues.
Truly comprehending its systemic nature depends on ST, described as a fundamental element of learning and teaching science~\cite{hmelor2007fish} and comprising skills for interpreting complex, dynamic systems~\cite{verhoeff2018theoretical}. 
Embedding ST within STEM teaching and learning can help pupils recognize complex, interconnected processes and build the skills needed to approach sustainability issues from an interdisciplinary perspective.

Developing ST about Earth systems progresses through increasingly complex cognitive stages, described as a three-level hierarchy of skills~\cite{assaraf2005development}.
At the \textit{Analysis} level, pupils identify a system's key elements and the processes operating within it.
The \textit{Synthesis} level involves connecting and organizing these elements, recognizing their dynamic interactions, and grasping the system's cyclical characteristics.
The \textit{Implementation} level pertains to interpreting hidden system features, generalizing about how the system functions, and applying temporal reasoning to interpret past events and anticipate future outcomes.

Kim described ST through ``the iceberg''~\cite{kim1999introduction}, in which visible \textit{Events} (the tip) arise from deeper \textit{Patterns} (middle) and \textit{Systemic Structures} (base) built on underlying \textit{Mental Models}, emphasizing that ST involves recognizing interconnections and feedback arising from deeper patterns and structures. 
Monat and Gannon~\cite{monat2015what} extend this model, situating it within a broader, integrated ST framework and stressing that ST for learning requires perceiving and interpreting interrelationships, non-linearity, feedback, and emergence.



% %%%%%%%%%%%%%%%%%%%%%%%%%%%%%%%%%%%%%%%%%%%%%%%%%%%%%%%%%%%% %
\subsection{Systems Thinking in Tracing Carbon}

The carbon cycle is challenging for pupils, who often perceive only a few components, struggle to interpret carbon transformations, treat sub-cycle processes as isolated, and fail to relate them to everyday contexts. 
These challenges make it a particularly relevant domain for studying how ST skills emerge and develop through interaction.

Within TC, ST provides a conceptual lens for both task design and interpretation: the carbon-cycle task assesses not only whether pupils can identify individual reservoirs and carbon flows, but also whether they can connect them into a broader, more dynamic understanding of the system. 
Rather than treating interaction logs as a record of isolated actions, we analyze them as traces through which events, patterns, and higher-level structures related to ST can be examined, consistent with theory-informed learning analytics that couples low-level traces with higher-order reasoning constructs.
The task we consider in TC targets a holistic understanding of the global carbon cycle connected to the following ST skills:

\begin{itemize}[nosep]
  \item
  Recognition of carbon cycle components, such as reservoirs and transfer/transforming processes (Analysis).
  
  \item 
  Identification of the relationships between components of the terrestrial and ocean carbon sub-cycles and the cyclic nature of the system (Synthesis).
  
  \item 
  Understanding of hidden dimensions of the carbon cycle and generalization from the terrestrial carbon sub-cycle to the ocean carbon sub-cycle (Implementation).
\end{itemize}





%%%%%%%%%%%%%%%%%%%%%%%%%%%%%%%%%%%%%%%%%%%%%%%%%%%%%%%%%%%%%%%%
% RELATED WORK
%%%%%%%%%%%%%%%%%%%%%%%%%%%%%%%%%%%%%%%%%%%%%%%%%%%%%%%%%%%%%%%%

\section{Related Work}
\label{sec:related_work}
To situate our work, we review three interconnected areas: (1) learning analytics; (2) event sequence visualization; and (3) motif discovery and matching.

% %%%%%%%%%%%%%%%%%%%%%%%%%%%%%%%%%%%%%%%%%%%%%%%%%%%%%%%%%%%% %
\subsection{Visual Learning Analytics}

Learning analytics (LA) aims to understand and optimize learning processes and environments~\cite{siemens2013learning, selwyn2020reimagining}.
While the field has evolved from descriptive activity analyses toward predictive models, adaptive feedback, and learning dashboards~\cite{matcha2019analytics, schwendimann2017perceiving, masiello2024current, paulsen2024learning}, the growing volume and complexity of learner interaction data highlight fundamental limitations in conventional LA systems, motivating research on visual analytics and learning dashboards~\cite{chatti2012reference, vieira2018visual, clow2013overview}.
A central challenge is the persistent disconnect between data collection and actionable pedagogical responses~\cite{clow2012learning}. Even when LA systems deliver predictions or summaries, educators' ability to translate them into instructional interventions remains limited~\cite{gasevi2015lets, jivet2018license}.

This is especially problematic in complex learning tasks, where aggregate statistics reveal little about \textit{how} learners approach problems, as reflected in the sequential and temporal organization of their behavior~\cite{vieira2018visual, verbert2013learning}. 
Without these behavioral traces, distinguishing deliberate strategies from trial-and-error behavior, and relating observable interactions to underlying reasoning processes becomes challenging~\cite{roll2015understanding}.
VA addresses this by supporting educators in formulating and testing hypotheses about student understanding~\cite{keim2008visual}, enabling exploration of behavioral processes in interpretable and pedagogically meaningful ways, as exemplified by recent interactive, learner-centered dashboards~\cite{gomez2021applying, khosravi2021intelligent, mohseni2024visual}.
However, support for data mining (DM) remains relatively uncommon in such tools~\cite{mohseni2024visual}. 
We extend this line of work by integrating visual analysis with DM, including motif discovery and clustering, within an interface grounded in systems thinking.



% %%%%%%%%%%%%%%%%%%%%%%%%%%%%%%%%%%%%%%%%%%%%%%%%%%%%%%%%%%%% %
\subsection{Learner Interaction Sequence Visualization}

Visual representations of interaction sequences enable richer sense-making of behavioral patterns than statistical summaries alone~\cite{chang2009defining, kang2009evaluating}, with chart- and timeline-based approaches summarized in recent surveys~\cite{yeshchenko2024survey, guo2022survey, eltayeby2016survey}.
TimeArcs, for example, track how pairwise relationships between entities evolve in dynamic networks~\cite{dang2016timearcs}. Emphasis on pairwise connectivity makes it less suitable for our setting, where the task is to preserve and compare learners' ordered interaction trajectories.
Marey-charts, originally time-space diagrams for train schedules~\cite{marey1885la, tufte1983visual}, visualize entities as lines with time on one axis and ordered states on the other, compactly capturing progression, pace, and transitions, and have since been adapted to workflow and process analytics~\cite{denisov2018unbiased, elzen2025towards, rubensson2025timeline}.
We adopt this design to represent learner progression, letting each path's shape reveal salient problem-solving behaviors and patterns.

In educational contexts, sequence visualizations increasingly target learner interaction logs from MOOCs, online platforms, and problem-solving environments, supporting exploration through sequence analysis, clustering, and multi-level comparison.
Examples include recovering learner trajectories from aggregate MOOC data~\cite{chen2020viseq}, identifying problem-solving logic and mistakes~\cite{supianto2016visualizations, xia2021qlens, tsung2022blocklens}, mapping interactions to cognitive processes~\cite{reda2014evaluating, yu2025revealing}, and deriving high-level sense-making from low-level events~\cite{gotz2009characterizing}.
Theory-informed analytics are grounded in structured frameworks from education, psychology, and cognitive science~\cite{hillaire2016prototyping}. Prior work has coupled log data with collaborative learning frameworks~\cite{xing2015group}, and revealed learners' cognitive states~\cite{chou2017open, mejia2017novel, pan2022theory}.
Yet, such couplings have primarily targeted pedagogical phases or task progressions rather than the underlying reasoning levels those actions reflect.
While recent work has organized learning environments around an ST hierarchy~\cite{mani2026designing, julsgard2026supporting}, the visual analysis of learners' sequential interactions and their relation to ST levels remains underexplored.
Because ST further spans distinct reasoning modes, from interconnectedness and feedback loops to synthesis and emergence~\cite{cabrera2019what}, visualizations should accommodate multiple levels of interpretation.
We therefore couple interaction sequence visualizations with ST reasoning hierarchies, enabling analysts to hypothesize connections between fine-grained behaviors and ST levels.



% %%%%%%%%%%%%%%%%%%%%%%%%%%%%%%%%%%%%%%%%%%%%%%%%%%%%%%%%%%%% %
\subsection{Pattern Discovery in Temporal Interaction Data}

Motif discovery finds recurring subsequences and supports the summarization and exploration of temporal datasets, including DNA sequences, motion data, and event logs~\cite{lin2002finding, liu2005locating, keogh2005clustering, fu2011review, baker2016educational, guerrini2025time}.
We adopt this perspective on motifs as recurring behavioral patterns and extend it to multidimensional representations derived from interaction data.
Sequence mining methods, such as PrefixSpan and SPADE, identify frequent ordered subsequences in event logs~\cite{pei2001prefixspan, zaki2001spade}, and have been used to analyze the evolution of learners' behaviors in open-ended environments~\cite{kinnebrew2014analyzing, zhang2023sequential, verstege2024using}.
These methods typically operate on discrete, categorical events and rely on exact or near-exact recurrence, limiting flexibility under temporal variation.
Process mining~\cite{romero2010handbook} has similarly been applied to educational data to model learner pathways, yet the resulting representations often lack the interactivity and multi-scale exploration needed for nuanced behavioral analysis~\cite{juhanak2019using, yu2026visual}.

Temporal networks align well with our data because interaction logs can be modeled as time-stamped relations (edges) between interface elements (nodes)~\cite{moody2005dynamic}, and temporal motifs~\cite{kovanen2011temporal, kovanen2013temporal} capture recurring time-respecting patterns as edges between entities over time~\cite{paranjape2017motifs, chen2025discovery}.
However, this approach poses several challenges for our setting.
Beyond the large edge count and computational cost, temporal network representations vary in how they encode event duration (e.g., instantaneous versus interval-based)~\cite{saryce2025powerful}, whereas our interactions vary substantially in duration.
Additionally, many also rely on snapshot-based partitioning into fixed intervals, which can obscure short-lived patterns and impose arbitrary boundaries on learner actions~\cite{liu2023temporal}.
Finally, temporal network motifs emphasize strict event relations and ordering~\cite{paranjape2017motifs}, whereas in our setting variations in action order (e.g., reversed sequences) do not necessarily reflect different strategies.

In time-series mining, matrix profile approaches identify motifs as pairs of approximately recurring patterns within a single series~\cite{yeh2016matrix, zhu2016matrix}, and a related extension considers consensus motifs across a set of series, broadening applicability to complex behavioral streams~\cite{kamgar2019matrix}.
This is especially relevant to our work, where multidimensional time series are derived from interaction sequences using a sliding window, and the problem is transformed into motif discovery over multidimensional series~\cite{yeh2017matrix}.
We therefore build on these lines by combining consensus and multidimensional motif discovery.




%%%%%%%%%%%%%%%%%%%%%%%%%%%%%%%%%%%%%%%%%%%%%%%%%%%%%%%%%%%%%%%%
% VA DESIGN REQ. AND TASKS
%%%%%%%%%%%%%%%%%%%%%%%%%%%%%%%%%%%%%%%%%%%%%%%%%%%%%%%%%%%%%%%%

\section{Design Process and Requirements}
\label{sec:design_process}
We describe the collaborative design process underlying \textsc{BEHAVE} and the requirements that informed its development.

% %%%%%%%%%%%%%%%%%%%%%%%%%%%%%%%%%%%%%%%%%%%%%%%%%%%%%%%%%%%% %
\subsection{Collaborative Design Process}
\label{subsec:design_proc}
\textsc{BEHAVE} was developed through a collaborative user-centered design and development process~\cite{christmann2022visualizing} involving an interdisciplinary team of co-authors of this paper.
The team comprised domain experts spanning three roles: visual learning and communication (VLC), design of the TC environment, and TC development.
The experts are also members of the TC research group, with expertise spanning STEM education, ST, and learning about the carbon cycle as a system. This same group designed TC, implemented its embedded learning tasks and logging architecture~\cite{mani2026designing}, and subsequently published an empirical study of pupils' engagement and strategy use with the environment~\cite{host2026pupils}.
Their dual role as both the developers of TC and analysts of its data gave them deep pedagogical and domain knowledge that directly shaped the requirements, behavioral abstractions, and the interpretation of observed patterns throughout the design process.

We engaged in iterative prototyping and workshops to frame the analytical problem, elicit requirements, and refine the design.
The first workshops established a shared terminology, familiarized us with the TC environment and its interaction log data, and identified the main analytical challenges faced by the experts.
These sessions clarified that existing approaches, based on aggregate statistics and manual inspection of selected logs, were insufficient for understanding the strategic and complex aspects of pupil behavior, motivating our design toward sequence-level representations (see also~\cref{subsec:va_system_design}).
From these discussions, we derived the core analytical needs that subsequently informed the design requirements (R$1$--R$4$) and tasks (T$1$--T$9$) presented in~\cref{subsec:va_system_design}.

We then entered an iterative co-development phase in which progressively more complex prototypes were presented and discussed, allowing experts to reflect on design choices, suggest revisions, and generate new analytical questions prompted by the evolving system. 
Expert feedback shaped the design through each iteration and is detailed in~\cref{sec:va_system_overview}.
Once a sufficiently complete prototype was established, we conducted final exploratory sessions in which domain experts examined the data in depth, assessing analytical utility while identifying remaining limitations and opportunities for refinement. 
The design process thus moved from early-stage problem framing, through iterative co-development, to situated exploratory use.



% %%%%%%%%%%%%%%%%%%%%%%%%%%%%%%%%%%%%%%%%%%%%%%%%%%%%%%%%%%%% %
\subsection{Systems Thinking as a Design Lens}

We adopt ST hierarchical levels~\cite{assaraf2005development} as a design lens to structure both the analytic workflow and the visual representations of \textsc{BEHAVE}. 
Our goal is to explore links between interaction logs and ST hierarchies. 
To this end, we use ST as a principal conceptual framework, beyond its implications for understanding Earth systems and for mapping pupil behavior in a digital problem-solving environment. 

ST emphasizes understanding a complex system such as the carbon cycle not only through its constituent parts, but also through the relationships among those parts, their evolution over time, and the higher-order patterns that emerge from their interaction. 
This aligns well with our data and analytic goals: pupils do not act through isolated events, but through temporally unfolding sequences whose meaning arises from transitions, recurrence, and progression across multiple interaction dimensions. 
Accordingly, we structure the interpretation of interaction behavior as a hierarchical progression of events, from identifying components, to recognizing relationships, dynamic processes, and cyclic behavior, to organizing these observations into a coherent interpretive framework that supports generalization, attention to hidden dimensions, and understanding behavior over time. 
Moreover, we link this hierarchical structure to ``the iceberg'' proposed by Kim ~\cite{kim1999introduction}, where components (Events) reflect the tip of the iceberg, relationships between components and pattern development (Patterns) correspond to the middle, and dynamic reasoning processes (Systemic Structures) and mental model construction represent the base. 
This framing helps design \textsc{BEHAVE} to support not only the detection of patterns but also their linking to ST reasoning.
In what follows, we make this derivation explicit.
Each requirement is introduced together with the analytical need that arose during the design process.


%%%%%%%%%%%%%%%%%%%%%%%%%%%%%%%%%%%%%%%%%%%%%%%%%%%%%%%%%%%%%%%%
\subsection{\textsc{BEHAVE} Design Requirements and Tasks}
\label{subsec:va_system_design}

The design process resulted in four high-level requirements (\textbf{R}$\mathbf{1}$--$\mathbf{4}$) that guide \textsc{BEHAVE}'s design and nine key analytical tasks (\textbf{T}$\mathbf{1}$-$\mathbf{9}$) that it should support.

% R 1 
\noindent
\textbf{R1: Representation of ST-related events.}
Since the experts found that aggregate statistics revealed \textit{that} particular behaviors occurred but not \textit{when} or \textit{in what order} they unfolded (\cref{sec:design_process}), \textsc{BEHAVE} should establish an interpretable visual mapping between logged interaction data and ST-related events. 
Derived from interaction logs through a domain-informed mapping that encodes behaviors relevant to ST, these events serve as analytically meaningful building blocks. 
At this level, the goal is to reveal how behavior unfolds over time and how the events relate to relevant attributes of the interaction.

\begin{enumerate}[label=T\arabic*., left=0cm, noitemsep, topsep=0pt, parsep=0pt]
    % T 1
    \item \textit{ST-related interaction event derivation and visualization}. 
    \textsc{BEHAVE} should visualize the ordering, interleaving, and duration of ST-related interactions over time.

    % T 2
    \item \textit{ST-related event characteristics inspection}. 
    \textsc{BEHAVE} should support the inspection of ST-related event dimensions, distribution, and contextual characteristics.
\end{enumerate}

% R 2
\noindent
\textbf{R2: Representation and exploration of ST-related patterns.}
Following the experts' observation that strategically meaningful behavior manifests as recurring combinations of events rather than isolated actions (\cref{sec:design_process}), \textsc{BEHAVE} should support identification and exploration of patterns understood as characteristic interrelationships among ST-related events. 
These ST-related patterns are important because they allow recurring strategic behavior to be studied as meaningful temporal patterns rather than as disconnected events. 
\textsc{BEHAVE} should therefore assist users in detecting repeated co-occurring ST-related events and examining how these combinations manifest within and across pupils, irrespective of their exact order.

\begin{enumerate}[label=T\arabic*., start=3, left=0cm, noitemsep, topsep=0pt, parsep=0pt]
    % T 3
    \item \textit{Motifs discovery}. 
    \textsc{BEHAVE} should provide a library of precomputed characteristic interaction patterns that may reflect distinct ST-related behaviors for users to inspect.
    
    % T 4
    \item \textit{Motifs visualization and analysis}. 
    \textsc{BEHAVE} should support exploration of how ST-related patterns are distributed across pupils and within complete interaction sequences.
\end{enumerate}

% R 3
\noindent
\textbf{R3: Explore underlying ST-related structures.}
Systemic structures indicate underlying patterns of organization in pupil behavior that reflect how pupils understand and engage with TC.
These structures are not directly observable in the raw interaction logs. 
Nevertheless, they may become visible when pupils are grouped according to how ST-related and recurring motifs unfold over time. 
Such grouping supports the synthesis of cohort-level behavioral profiles and can provide initial indications of ST-related structures that help explain why different behavioral patterns emerge.
Because experts sought to understand \textit{why} such patterns emerge, and because this can only be inferred through grouping rather than observed directly (\cref{sec:design_process}), \textsc{BEHAVE} should support exploration of such systemic structures.

\begin{enumerate}[label=T\arabic*., start=5, left=0cm, noitemsep, topsep=0pt, parsep=0pt]
    % T 5
    \item \textit{Cohort identification}. 
    \textsc{BEHAVE} should support the detection of clusters with similar behavioral profiles based on survey variables, behavioral descriptors, and sequence-level statistics.

    % T 6
    \item \textit{Cohort characterization and inspection}. 
    For each identified cohort, \textsc{BEHAVE} should support interpretation of the basis of similarity through the distribution of descriptive features and characteristic behavioral tendencies.
    This should be coupled with a detailed, individual inspection so that users can examine the overarching workflow of an individual of interest.

    % T 7
    \item \textit{Inter-cohort comparison}. 
    \textsc{BEHAVE} should support the comparison of characteristic ST-related patterns within and across cohorts to identify possible indications of underlying ST-related structures.
\end{enumerate}

% R 4
\noindent
\textbf{R4: Human in the analysis loop.}
Because the design sessions established that relating observed behavior to ST is interpretive rather than purely algorithmic, requiring computational results to remain open to expert revision (\cref{sec:design_process}), \textsc{BEHAVE} should keep users actively involved throughout the workflow. 
Computational methods can suggest clusters and motifs, but these results must remain open to validation, reinterpretation, and refinement through domain expertise.
\textsc{BEHAVE} should therefore support flexible hypothesis exploration and interactive adjustment of computational outcomes by providing coordinated views.

\begin{enumerate}[label=T\arabic*., start=8, left=0cm, noitemsep, topsep=0pt, parsep=0pt]
    % T 8
    \item \textit{Coordinated views and interactivity}. 
    \textsc{BEHAVE} should include coordinated views and provide interactions to aid users in inspecting, comparing, validating, and refining findings, including investigation and interpretation of similarity, motifs, and cohorts.

    % T 9
    \item \textit{Hypothesis-driven exploration}. 
    \textsc{BEHAVE} should support the exploration of hypotheses about the relationships between ST-related events, motifs, cohorts, and ST reasoning.
\end{enumerate}

\begin{figure}
    \centering
    \includegraphics[
    width=.94\columnwidth,
    alt={}
    ]{figs/pipeline.pdf}
    \caption{
    Overview of \textsc{BEHAVE}'s analysis pipeline.
    Raw data are preprocessed into pupil logs and abstracted into multidimensional series, which serve as the main inputs to \textsc{BEHAVE}, consisting of three coordinated panels for clustering, motif investigation, sequence comparison, and individual inspection.
    }
    \label{fig:pipeline}
\end{figure}





%%%%%%%%%%%%%%%%%%%%%%%%%%%%%%%%%%%%%%%%%%%%%%%%%%%%%%%%%%%%%%%%
% SYSTEM DESCRIPTION
%%%%%%%%%%%%%%%%%%%%%%%%%%%%%%%%%%%%%%%%%%%%%%%%%%%%%%%%%%%%%%%%

\section{\textsc{BEHAVE} Overview and Description}
\label{sec:va_system_overview}

Our visual analytics pipeline comprises three stages, as illustrated in~\cref{fig:pipeline}.
\textsc{BEHAVE} corresponds to the final stage of this pipeline and consists of three primary components: the \textsc{Cluster} (\cref{subsec:cluster_panel}), \textsc{Comparison} (\cref{subsec:comparison_panel}), and \textsc{Individual} (\cref{subsec:individual_panel}) \textsc{Panels}. 
Together, these components support interactive exploratory analysis of pupils' problem-solving behavior logged in the TC environment through coordinated views and associated control panels. 
Each view includes a dedicated control panel that allows users to adjust visual encodings, select data points of interest, and apply filters, enabling exploration of data from an ST perspective. 

The interaction logs that \textsc{BEHAVE} analyzes are large, heterogeneous, and of variable-length, with meaning residing not in individual actions but in their ordering, recurrence, and co-occurrence across behavioral dimensions over time.
This places them in the same broad class as other temporally variable, multidimensional sequence problems (e.g., motion capture, event logs, sensor streams), making the pipeline below relevant beyond the specific Tracing Carbon environment.


%%%%%%%%%%%%%%%%%%%%%%%%%%%%%%%%%%%%%%%%%%%%%%%%%%%%%%%%%%%%%%%%
\subsection{Data Characterization and Preprocessing}
\label{subsec:data_characterization}

In TC, pupils construct carbon cycles by drawing directed arrows connecting carbon reservoirs.
Interactions are captured through fine-grained event logging that records the source and target reservoirs, the correctness of each arrow, and millisecond-resolution timestamps of interaction start and end. 

In this work, we illustrate our approach by applying it to the most complex task in TC and extracting the corresponding subset of interaction data from 97 pupils.
Each logged interaction contains a \textit{source} and a \textit{target}, indicating the starting and ending reservoir of an arrow, respectively. 
After an initial preprocessing step, as described in detail in the supplementary materials, each completed drag is represented by its start and end timestamps, from which we derive the duration of the action and the idle interval to the next interaction. 
We further grouped reservoirs into five semantically meaningful categories: \textit{atmosphere}, \textit{factory}, \textit{surface}, \textit{fossil fuel}, and \textit{ocean}, as illustrated in~\cref{fig:tcInterface}.

Additionally, building on prior manual analysis and discussions with domain experts, we identified five types of behavioral building blocks corresponding to elementary ST-related events:
\begin{itemize}[nosep]
  \item 
  \textbf{\textit{Tracing}}: Implies the action of connecting reservoirs successively, i.e., the current arrow's end point is the next's starting point.
  
  \item 
  \textbf{\textit{Minimal sub-cycle}} (hereafter \textbf{\textit{mini-cycle}}): Corresponds to creating a cycle between pairs of reservoirs, i.e., drawing an arrow from one reservoir to another and back again.

  \item 
  \textbf{\textit{Same-start}}: Corresponds to repeatedly drawing arrows from one reservoir to several different reservoirs.
  
  \item 
  \textbf{\textit{Same-end}}: Corresponds to repeatedly drawing arrows from multiple reservoirs toward the same reservoir.
  
  \item 
  \textbf{\textit{Within-category}}: Implies making connections within the same category, i.e., the starting and ending reservoirs of a connection belong to the same category.
\end{itemize}

In addition to these five ST-related events, we extracted two measures that further characterize pupils' interaction behavior.
\begin{itemize}[nosep]
  \item
  \textbf{Sparsity}: A temporal measure of how active a pupil is over a given period, such that high sparsity corresponds to longer idle intervals between actions.

  \item 
  \textbf{Correctness}: A Boolean measure indicating whether the connection made was correct. 
\end{itemize}

Based on these descriptors, we computed corresponding binary values for each interaction sequence. 
For example, if two or more consecutive dragging interactions originated from the same reservoir, a value of $1$ is assigned for \textit{same-start}.
This procedure was applied to all five ST-related events and the correctness measure, yielding six features per sequence.


%%%%%%%%%%%%%%%%%%%%%%%%%%%%%%%%%%%%%%%%%%%%%%%%%%%%%%%%%%%%%%%%
\subsection{Data Abstraction and Processing}
\label{subsec:data_abstraction}

The input to \textsc{BEHAVE} is a collection of pupil interaction logs, each containing discrete actions with timestamps, six binary feature vectors, and sparsity as described in~\cref{subsec:data_characterization}.
Our pipeline transforms these logs into multidimensional time series and computes motifs.
The pipeline involves a small set of parameters in two classes, summarized in \cref{tab:parameters}.
\textit{Data-derived} parameters are computed from dataset statistics, while \textit{task-related} parameters are set depending on the analysis focus to balance sensitivity to recurring patterns against robustness to noise during computation. 
Our reasoning for setting the task-related parameters is reported below and detailed further in the supplementary materials. 
Although changing the \textit{task-related} parameters may alter the resulting set of discovered motifs, the findings reported in the case study (\cref{sec:case_study}) and validation (\cref{sec:expert_validation}) did not appear sensitive to the exact value of any single parameter. 
However, a more systematic sensitivity analysis remains part of future work (\cref{sec:discussion}).

\begin{table*}[t]
\centering
\setlength{\tabcolsep}{18pt}  
\caption{Parameters of the \textsc{BEHAVE} pipeline. \textit{Data-derived} parameters are computed from dataset statistics; \textit{task-related} parameters are set to balance match and merge granularity during motif discovery.
Derivations of the values are reported in the supplementary materials.}
\label{tab:parameters}
\begin{tabular}{@{}llll@{}}
\toprule
Parameter & Role & Basis & Value \\
\midrule
\multicolumn{4}{@{}l}{\textit{Data-derived}} \\
$W$ & Sliding-window width & Mean active-segment duration & $15$\,s \\
$g$ & Active-segment split & Mean inter-interaction gap & $\approx 4$\,s \\
$\rho$ & Sparsity calibration & $5$th percentile of inter-event gaps & $1.448\,s$ \\
$m_{\min}$ & Minimum support & Lower quartile of action counts & $5$ \\
\midrule
\multicolumn{4}{@{}l}{\textit{Task-related}} \\
$\theta_{\mathrm{pool}}$ & Activity-pool cutoff & Conservative activity gate & $0.1$ \\
$\tau$ & Match-distance threshold & Sensitivity vs.\ noise robustness & $0.12$ \\
$\varepsilon$ & Merge tolerance & Near-duplicate aggregation ($\pm5\%$) & $0.05$ \\
\bottomrule
\end{tabular}
\end{table*}



%%%%%%%%%%%%%%%%%%%%%%%%%%%%%%%%%%%%%%%%%%%%%%%%%%%%%%%%%%%%%%%%
% \subsubsection{Multidimensional Series Construction}
\noindent
\textbf{Multidimensional series construction}.
\label{subsubsec:multi_series}
To capture local temporal patterns in pupils' interaction behavior, we summarize each sequence as a multidimensional time series derived from overlapping sliding windows, so that each feature forms a signal reflecting the pupil's local behavioral tendency over the course of the interaction.
A sliding window of width $W{=}15$\,s advances in one-second steps across each sequence.
To determine an appropriate temporal threshold for separating distinct periods of activity, we first computed an average gap duration between consecutive interactions across all interaction sequences ($g\approx4$\,s). 
Using this threshold, each interaction sequence was partitioned into \emph{active segments}, such that a new segment was initiated whenever the gap between two consecutive interactions exceeded $4$\,s. 
Following this, we computed the mean active-segment duration for the sequence ($=14.285$\,s), rounded to $15$\,s, and used it as the sliding window width ($W$).


At each window position, only actions that both start and end within its bounds are considered. 
This yields six behavioral descriptors, each defined as a bounded ratio in $[0,1]$, computed as follows:

\begin{itemize}[nosep]
    \item
    \textit{Tracing} and \textit{mini-cycle} are each computed as the fraction of interactions within the window flagged as that behavior, and therefore lie in $[0,1]$.
    
    \item
    \textit{Same-start}, \textit{same-end}, and \textit{within-category} are each measured from the longest consecutive corresponding behavior within a window.
    \begin{equation}
        \operatorname{rate(w_s)}=
        \begin{cases}
        \dfrac{R_{\max}(W_s)}{n}, & \text{if } R_{\max}(W_s)\ge 2 \text{ and } n>0, \\[0.6em]
        0, & \text{otherwise,}
        \end{cases}
    \end{equation}
    where $W_s$ is the start time of a sliding window, $R_{\max}(W_s)$ is the length of the longest occurrence of the feature, and $R_{\max}\ge 2$ assigns a nonzero rate only when the behavior repeats. 
    Since a run cannot exceed the number of actions, this rate bounds in $[0,1]$.
    A higher value indicates a longer consecutive streak of the behavior within the window.
    
    \item
    \textit{Sparsity} reflects how active a pupil is during a window.
    \begin{equation}
        \mathrm{sparsity} = 1 - \frac{n_{\mathrm{actions}}}{W / \rho}
    \end{equation}
    where $n_{\mathrm{actions}}$ is the number of actions in the window $W$, and $\rho$ is a data-derived reference inter-action time.
    A higher value indicates long idle time between few actions, and vice versa.
\end{itemize}

Correctness is excluded from the feature vector used for motif discovery, since we aim to identify recurring problem-solving strategies rather than outcomes.
It is retained separately to support later analysis of associations between strategies and answer accuracy.



%%%%%%%%%%%%%%%%%%%%%%%%%%%%%%%%%%%%%%%%%%%%%%%%%%%%%%%%%%%%%%%%
\noindent
\textbf{Motif discovery}.
Based on the derived multidimensional time series, we perform motif discovery inspired by consensus~\cite{kamgar2019matrix} and multidimensional motifs~\cite{yeh2017matrix} to identify recurring local behavioral patterns across pupils.

We define $\mathbf{x}_{s,t} \in \mathbb{R}^D$ as the descriptor vector of index $t$ in sequence $s$, where the $D=6$ dimensions correspond to the six interaction behaviors introduced in~\cref{subsec:data_characterization} (e.g., $t_1$ contains the window spanning from 1\,s to 15\,s, and $t_2$ contains the window spanning from 2\,s to 16\,s, etc.).
All window vectors are pooled across sequences, and windows with negligible activity are excluded before matching via the activity-pool cutoff $\theta_{\mathrm{pool}}$, which acts as a conservative pre-filter separating behaviorally meaningful windows from near-inactive ones.
Because it discriminates activity from inactivity rather than one behavior from another, a low cutoff suffices, and the downstream motifs are insensitive to its exact value.
Windows near the threshold contain little activity and contribute few actions to any candidate motif.

Two windows are considered a match when the distance between their descriptor vectors falls below the match-distance threshold $\tau$:
\begin{equation}
d\bigl(\mathbf{x}_{s,t}, \mathbf{x}_{s',t'}\bigr) < \tau
\end{equation}
where $d(\cdot,\cdot)$ is the Euclidean distance between two descriptor vectors.
If $\tau$ is set too low, the threshold admits only near-identical windows, which risks filtering out meaningful recurring patterns; if set too high, it may merge behaviorally distinct windows, resulting in many and inconsistent motifs.
For each candidate window in the pool, $(s,t) \in \mathcal{P}$, we collect all matching windows from other pupils, and rank candidates by their cross-sequence support, defined as the number of distinct pupils containing at least one match.
This favors motifs recurring across pupils rather than repeatedly within a single sequence.

Subsequently, temporally adjacent candidates within the same sequence are consolidated into a single representative motif to reduce redundancy, and matches are deduplicated so that no two retained matches overlap within the same sequence.
We discard candidates below the \textit{data-derived} minimum interaction count $m_{\min}$ to avoid sparse-window artifacts, and merge motifs whose feature values are near-identical:
\begin{equation}
|p_j - q_j| \le \varepsilon,
\qquad
\forall j \in \mathcal{J}
\end{equation}
where $\mathbf{p}$ and $\mathbf{q}$ are two candidate motifs, $\varepsilon$ is the near-duplicate merge tolerance, and $\mathcal{J}$ is the set of compared feature dimensions.
A larger $\varepsilon$ yields a smaller, coarser motif library, and a smaller $\varepsilon$ results in a larger, more redundant one.
After merging, the match lists are unioned, the per-pupil non-overlap constraint is reapplied, and motif support is recomputed and finalized.

As a result, the discovered motifs represent recurring, temporally localized, and non-redundant behavioral patterns shared across pupils, which serve as inputs to \textsc{BEHAVE} for further exploration.
Additional details on the parameters are provided in the supplementary materials.



%%%%%%%%%%%%%%%%%%%%%%%%%%%%%%%%%%%%%%%%%%%%%%%%%%%%%%%%%%%%%%%%
\subsection{Individual Panel}
\label{subsec:individual_panel}

The \textsc{Individual Panel} provides the most fine-grained level of analysis, supporting detailed inspection of a single pupil's interaction trajectory. 
It operationalizes \textbf{R1} by preserving and visualizing the sequential structure of ST-related events, allowing users to inspect how behavior unfolds over time, including the ordering, interleaving, and duration of interactions and derived ST-related events (\textbf{T1}).

We adopt a Marey-chart-based representation of interaction events (\cref{fig:teaser}) with two complementary modes, a design that emerged from the co-design process: when we initially discarded duration to normalize sequences, the domain experts countered that pacing was itself analytically meaningful, prompting us to preserve both perspectives.
In \textit{sequential mode}, all stringlines are assigned equal length (\cref{fig:trial_error}.2), suppressing real-time duration in favor of interaction order.
This makes the view effective for identifying task-solving strategies, repeated local patterns, shifts in behavior (\textbf{T2}), and for inspecting computed motifs in their sequential context to assess whether algorithmically derived patterns are behaviorally meaningful (\textbf{T4}).
In \textit{durational mode}, each stringline is scaled to the actual interaction duration (\cref{fig:trial_error}.3), reflecting both drawing time and idle intervals, which is useful for examining temporal progression, pacing, and whether activity occurs in bursts or with an even rhythm (\textbf{T2}).
Because pupils often create connections in rapid succession, this view can become dense. We therefore adopt a comet-path convention to distinguish each interaction's start and end points~\cite{holten2011extended, armstrong2014visualizing}.
In both modes, reservoirs are spatially grouped by category, a structure identified by the experts as central to TC's organization, enabling abstraction to the level of pupils' overall reasoning flow, while background colors provide semantic support, e.g., azure for atmosphere-related reservoirs.

Above the Marey-chart, the panel displays the multidimensional time series derived from the sliding-window abstraction (\cref{subsec:data_abstraction}), summarizing the temporal tendency of the six behavioral descriptors and allowing users to examine not only which behaviors occur but also how they evolve (\cref{fig:teaser}.3). 
The legend doubles as an interactive filter for displaying selected series on demand based on the analysis goal.

To further support inspection of ST-related event characteristics (\textbf{T2}), the panel includes the derived ST-related event sequence (\cref{subsec:data_characterization}), with color encoding correctness (\cref{fig:teaser}.5). 
This provides a compact overview of which events occur, when, and whether they are correct, letting users quickly identify periods dominated by particular tendencies and relate them to task outcomes.
Finally, because several main ST-related events are mutually exclusive, e.g., \textit{same-start} cannot co-occur with \textit{same-end} or \textit{tracing}, we bundle them into a compact single-row visualization to reduce redundancy and improve interpretability (\textbf{T1}) (\cref{fig:teaser}.6). 
\textit{Mini-cycle}, a special form of \textit{tracing}, shares its hue but is distinguished through higher saturation.



%%%%%%%%%%%%%%%%%%%%%%%%%%%%%%%%%%%%%%%%%%%%%%%%%%%%%%%%%%%%%%%%
\subsection{Motif and Comparison Panel}
\label{subsec:comparison_panel}

The \textsc{Motif and Comparison Panel} supports visualizing and comparing pupils' ST-related event sequences, enabling exploration of recurring temporal patterns and their distribution across pupils and cohorts. 
Users can thereby move beyond isolated events and study characteristic ST-related patterns as repeated subsequences embedded in full sequences.

Its central component is the \textit{Motif Library} (\cref{fig:teaser}.1), displaying discovered motifs (\textbf{T3}) that may correspond to distinct ST-related strategies or difficulties. 
Motifs are sorted by cross-sequence support (\cref{subsec:data_abstraction}) by default, but can be reordered by one or more behavioral metrics, such as \textit{same-start} or \textit{same-end}, to suit the analytical goal (\textbf{T9}). 
For cross-sequence comparison, the compact single-row visualization from the \textsc{Individual Panel} is reused (\cref{fig:teaser}.2), trading per-dimension detail for a compact view of locally dominant behavior.
Users can select motifs to inspect their distribution across full sequences and cohorts, e.g., \textit{where}, \textit{when}, and how frequently they occur and whether they concentrate in particular clusters (\textbf{T4}), and can group sequences by cluster to compare within and across cohorts (\textbf{T7, T8}).

This multi-sequence comparison allows users to determine whether a cluster reflects a genuinely shared behavioral structure or only coarse-grained similarity across sequences, bridging cohort-level summaries and single-sequence inspection. 
The panel thus supports iterative interpretation: comparing candidate explanations, assessing whether motifs are characteristic of particular cohorts, and forming new hypotheses about the relation between recurring patterns and ST reasoning (\textbf{T9}).



%%%%%%%%%%%%%%%%%%%%%%%%%%%%%%%%%%%%%%%%%%%%%%%%%%%%%%%%%%%%%%%%
\subsection{Cluster Panel}
\label{subsec:cluster_panel}

The \textsc{Cluster Panel} provides a cohort-level overview and supports \textbf{R3} by making potential underlying ST-related structures visible through clustering. 
It supports identifying cohorts with similar behavioral profiles (\textbf{T5}), characterizing them through visual summaries (\textbf{T6}), and comparing cohort attributes (\textbf{T7}).

Clustering operates over a flexible feature set (\textbf{T5}): survey variables (e.g., gender and pupils' task ratings of fun, instructive, or interesting, collected after each task), behavioral descriptors such as aggregate ST-related event rates (\cref{subsec:data_characterization}), and sequence-level statistics such as total duration, number of interactions, and correctness ratio.
This allows users to construct cohorts from behavioral, performance-related, or survey-based perspectives. 
Because features may be numerical or categorical, we employ \textit{k}-prototypes~\cite{huang1998extensions} (reducing to \textit{k}-means~\cite{mcqueen1967some} for purely numerical input), with an elbow-method helper suggesting an initial cluster count that users can adjust interactively.
Results are shown in a two-dimensional scatterplot, where each point represents a pupil colored by cluster membership, conveying the separation, compactness, and size of the clusters (\cref{fig:cluster_detail}.2). 
The axes can represent individual selected features directly or a $2$D projection of the feature space.

To support cohort characterization (\textbf{T6}), per-cluster summaries are displayed as donut charts for categorical attributes such as gender, and bar charts of cohort averages for numerical attributes (\cref{fig:cluster_detail}.1). 
These help users interpret the basis of similarity within a cohort, for example, clustering by ST-related events and then assessing whether cohorts differ in survey responses or accuracy. 
For inter-cohort comparison (\textbf{T7}), a feature-centered \textit{Comparison View} displays a selected attribute across clusters using aligned bar charts (\cref{fig:four_clusters}). 
This is important in our setting because systemic structure is not directly observable in the logs and must be inferred from the distribution of behaviors and feedback across pupils. 
The \textsc{Cluster Panel} thus serves as an entry point for examining cohort structure, interpreting how cohort-level differences may reflect underlying ST-related organization, and driving hypothesis-driven exploration (\textbf{T9}). 
This interactive, on-the-fly clustering emerged from the co-design process. 
Experts found precomputed community structures insufficient for their needs and preferred to define cohorts dynamically through feature selection and parameter adjustment.

\begin{figure}
    \centering
    \includegraphics[
    width=.91\columnwidth,
    alt={
    Four identified cohorts (C0–C3). 
    (1) Each cohort is summarized by a donut chart indicating cohort size and gender composition, alongside bar charts encoding survey-derived, aggregated behavioral, and sequence-level features. 
    (2) A two-dimensional scatterplot projects pupils onto the prototype-distance space, with point color encoding cohort membership.
    }
    ]{figs/cluster_detail.pdf}
    \caption{
    Four identified cohorts ($C0$-$C3$). 
    (1) Each cohort is summarized by a donut chart indicating cohort size and gender composition, alongside bar charts encoding survey-derived, aggregated behavioral, and sequence-level features. 
    (2) A two-dimensional scatterplot projects pupils onto the prototype-distance space, with point color encoding cohort membership.
    }
    \label{fig:cluster_detail}
\end{figure}



%%%%%%%%%%%%%%%%%%%%%%%%%%%%%%%%%%%%%%%%%%%%%%%%%%%%%%%%%%%%%%%%
\subsection{Interactions with \textnormal{\scshape BEHAVE}}
\label{subsec:interaction_with_va}

\textsc{BEHAVE} supports two complementary workflows with distinct entry points: a cluster-driven workflow that begins with cohort segmentation, and a motif-driven workflow that begins with frequently recurring patterns. 
Both workflows converge in the \textsc{Individual Panel} for fine-grained inspection (\textbf{T8}). The analytical process is iterative, with users navigating between panels as new observations prompt hypothesis refinement (\textbf{T9}).

\noindent
\textbf{Cluster-driven workflow}.
Users define a feature space and set the number of clusters directly or via the elbow-based helper. 
The resulting cohorts can be inspected through their summaries and compared on survey-derived, behavioral, or sequence-level features (\cref{fig:cluster_detail}, \cref{fig:four_clusters}) (\textbf{T5, T6, T7}). 
This identifies clusters that merit further investigation, such as cohorts with higher or lower correctness when grouped by ST-related events. 
Once a cohort has been identified, users can examine whether it exhibits distinctive behavioral patterns in the \textsc{Motif and Comparison Panel}.

\noindent
\textbf{Motif-driven workflow}.
Alternatively, the \textsc{Motif and Comparison Panel} supports exploration based on the prevalence of specific recurring patterns rather than cohort structure (\textbf{T3}).
The \textit{Motif Library}, sorted by cross-sequence support or by a selected ST-related event, lets users begin from either frequent or behavior-specific motifs. 
Once motifs are selected, sequences in the \textit{Comparison View} are reordered by match count, foregrounding those with stronger motif presence (\textbf{T4}); hovering over or clicking a motif isolates its occurrences and reveals the surrounding context.
This lets users assess whether a motif is concentrated in particular cohorts (\textbf{T7}) and form hypotheses about how recurring patterns relate to ST (\textbf{T9}).

\noindent
\textbf{Individual-level inspection}.
When a pupil becomes relevant for closer inspection, users select that sequence and transition to the \textsc{Individual Panel} to examine temporal progression, local tendencies, and the relationship between recurring motifs and the broader context. 
Abstract patterns identified at the cohort or motif level can thus be related to concrete behavioral evidence (\textbf{T6}), for example by examining whether a motif reflects a particular ST stage or a broader problem-solving tendency (\textbf{T1}, \textbf{T2}).





%%%%%%%%%%%%%%%%%%%%%%%%%%%%%%%%%%%%%%%%%%%%%%%%%%%%%%%%%%%%%%%%
% CASE STUDY
%%%%%%%%%%%%%%%%%%%%%%%%%%%%%%%%%%%%%%%%%%%%%%%%%%%%%%%%%%%%%%%%

\section{Case Study}
\label{sec:case_study}
In this study, we used \textsc{BEHAVE} to move from broad exploration toward the formulation and refinement of hypotheses concerning pupils' ST skills, which involve recognizing connections and cyclical processes.
The TC interactive digital environment supports this analysis by prompting pupils to draw arrows that represent carbon-transition relationships.

Rather than determining pupils' ST levels, which interaction logs alone cannot establish, the case study demonstrates how \textsc{BEHAVE} surfaces candidate behavioral indicators for further targeted investigation.
The analysis reported here represents the final stage of our co-design process (\cref{subsec:design_proc}), consisting of in-depth exploratory sessions in which the domain experts analyzed the TC interaction data using \textsc{BEHAVE}. 
Because these experts are also co-designers of the learning environment, they bring deep pedagogical and domain knowledge to the interpretation of behavioral patterns. 
Throughout the analysis, we treat the resulting interpretations as hypotheses to be investigated rather than validated conclusions, and we explicitly identify where external assessment data would be required for confirmation (\cref{sec:discussion}). 

Working with \textsc{BEHAVE}, the experts identified recurring interaction motifs and examined behavioral patterns across pupil groups clustered by multiple characteristics. 
These motifs supported the formulation of hypotheses about the relationship between behavioral patterns and task performance, for example, whether pupils who frequently employed motifs with strong cyclic and tracing patterns were also more accurate in connecting carbon cycle elements.
When examining pupils clustered by interaction correctness, the experts similarly formulated hypotheses linking higher correctness to increased cyclic and tracing behaviors, as well as to spatial and temporal characteristics such as sparsity or within-category arrow drawing.
Drawing on their pedagogical knowledge, the experts in turn used these observations to reason about pupils' learning processes, hypothesizing, for instance, that applying ST-related event patterns such as \textit{tracing} and \textit{mini-cycles} may promote deeper ST understanding, which might manifest in fewer errors and shorter idle times. 



%%%%%%%%%%%%%%%%%%%%%%%%%%%%%%%%%%%%%%%%%%%%%%%%%%%%%%%%%%%%%%%%
\subsection{Strategy Profiles}

To better understand variation in problem-solving approaches, the experts first clustered pupils based on the four previously defined strategies~\cite{host2026pupils}, with the number of clusters set to four following the elbow method.
These four clusters exhibit different compositions of ST-related event usage, indicating different strategies.
One group ($C2$) is distinguished by a relatively high proportion of \textit{tracing} behavior and \textit{mini-cycles}, suggesting a tendency to follow paths through the carbon cycle.
Another group ($C3$) is dominated by \textit{same-start} events, while a third group ($C0$) shows a high proportion of \textit{same-end} events.
The fourth group ($C1$) combines a high fraction of both \textit{same-start} and \textit{same-end}, as illustrated in~\cref{fig:four_clusters}.

These behavioral differences are reflected in performance outcomes. 
The group with a high proportion of \textit{same-end} strategy ($C0$) achieves the highest correctness scores.
In contrast, the group combining \textit{same-start} and \textit{same-end} strategies ($C1$) is associated with the lowest correctness. 
The experts interpreted this as an indication that not all systematic exploration strategies are equally effective, and that certain patterns of repetition may reflect less productive reasoning.

\begin{figure}
    \centering
    \includegraphics[
    width=\columnwidth,
    alt={
    A screenshot of four clusters characterized by different compositions of ST-related event usage, indicating different strategies.
    The bar-chart summaries show the relative levels of same-start, same-end, tracing, and mini-cycle rates for each cluster, providing a compact behavioral profile of the groups.
    }
    ]{figs/four_clusters.pdf}
    \caption{
    Different compositions of strategy usage characterize four clusters.
    The bar-chart summaries show the relative levels of \textit{same-start}, \textit{same-end}, \textit{tracing}, and \textit{mini-cycle} rates for each cluster. 
    }
    \label{fig:four_clusters}
\end{figure}

The distribution of the four strategies can be further explored by visualizing the occurrence of identified motifs across pupils in the different clusters. 
Many motifs were characterized by a single ST-related event, confirming that pupils often performed actions aligned with a single strategy at a time.
This can be followed up by investigating selected pupils in the \textsc{Individual Panel}, where distinct shifts in apparent strategy are often discernible. 
For example, as illustrated in~\cref{fig:switch_strategies}, the pupil starts by tracing a traditional food web, in which plants are consumed by animals, which subsequently die, and are decomposed by decomposers. 
Drawing on their knowledge of common carbon-cycle misconceptions, the experts noted that a further \textit{tracing} move connecting decomposers back to plants may reflect the well-known misconception that plants obtain most of their nutrients from the soil rather than through photosynthesis.
The pupil then performs further moves dominated by \textit{tracing}, before switching to the \textit{same-start} strategy to explore how ocean carbon dioxide and the atmosphere each connect to the rest of the carbon cycle.

\begin{figure}
    \centering
    \includegraphics[
    width=.91\columnwidth,
    alt={
    A screenshot of an individual pupil’s interaction sequence and local behavioral tendencies. 
    Early in the sequence, the pupil mainly follows a tracing strategy, exploring a traditional food-web path from plants to animals and decomposers, including a possible misconception-driven attempt to connect decomposers back to plants. 
    Later, the behavior shifts toward a same-start strategy, with repeated exploration from ocean carbon dioxide and the atmosphere to investigate their connections to the larger carbon cycle.
    }
    ]{figs/switch_strategies.pdf}
    \caption{
    Example of a pupil’s interaction sequence and behavioral tendencies. 
    Early in the sequence, the pupil mainly follows a \textit{tracing} strategy.
    Later, the behavior shifts toward a \textit{same-start} strategy.
    }
    \label{fig:switch_strategies}
\end{figure}

These observations deepen our understanding of the data and highlight that pupils adopt diverse problem-solving strategies, with the relationship between ST-related strategies and correctness appearing multifaceted rather than uniform.





%%%%%%%%%%%%%%%%%%%%%%%%%%%%%%%%%%%%%%%%%%%%%%%%%%%%%%%%%%%%%%%%
\subsection{Uncovering Trial-and-Error Behavior}

Beyond strategy composition, \textsc{BEHAVE} enables the analysis of pupils' use of trial and error in problem-solving. 
In the hypothesis that trial and error would be associated with less time spent deliberating over appropriate moves, the experts explored the data by clustering on \textit{sparsity} and, in turn, on the frequency of use of each ST-related event strategy.

One of the clusters based on \textit{tracing} and \textit{sparsity} was associated with a very low average sparsity value.
Inspecting these sequences in the \textsc{Motif and Comparison Panel}, the experts observed a peculiar pattern exhibited by several pupils, consisting of a consecutive sequence of \textit{mini-cycle} and \textit{tracing} moves, varying in length from three to $15$ consecutive moves.
Examining these pupils in the \textsc{Individual Panel}, as illustrated in~\cref{fig:trial_error}, revealed that the pattern consisted of repeatedly drawing mini-cycles back and forth from one reservoir to its surrounding reservoirs.
In conjunction with the low sparsity scores (i.e., little time between moves), the experts interpret this as a novel strategy that may be uniquely associated with trial-and-error problem-solving. 
To investigate this behavior in greater detail, they isolated a motif characterized by $50$ percent \textit{tracing} and $50$ percent \textit{mini-cycle} interactions.
They then explored its distribution across the entire dataset and found several additional pupils exhibiting the same pattern.
This suggests that the observed trial-and-error behavior was not an isolated case.
Overall, this analysis demonstrates how \textsc{BEHAVE} supports the discovery of previously uncharacterized behavioral patterns and the formulation of new hypotheses about pupils' problem-solving strategies.

\begin{figure}
    \centering
    \includegraphics[
    width=.95\columnwidth,
    alt={
    A screenshot of a trial-and-error behavioral pattern identified in \textsc{BEHAVE}: (1) the full interaction sequence, with the boxed segment shown enlarged in two complementary views, (2) in sequential mode, and (3) in duration mode. 
    Both zoom into the same highlighted segment of (1) and reveal repeated alternation between tracing and mini-cycle moves, often forming rapid back-and-forth connections from one reservoir to nearby reservoirs.
    Given the low sparsity of these sequences, this pattern suggests a distinct trial-and-error strategy that was later used as a manually defined motif for broader exploration of the dataset.
    }
    ]{figs/trial_error.pdf}
    \caption{
    Example of a trial-and-error behavioral pattern identified in \textsc{BEHAVE}: (1) the full interaction sequence, with the boxed segment shown enlarged in two complementary views, (2) in sequential mode, and (3) in duration mode. 
    Both zoom into the same highlighted segment of (1) and reveal repeated alternation between \textit{tracing} and \textit{mini-cycle} moves, often forming rapid back-and-forth connections from one reservoir to nearby reservoirs.
    }
    \label{fig:trial_error}
\end{figure}





%%%%%%%%%%%%%%%%%%%%%%%%%%%%%%%%%%%%%%%%%%%%%%%%%%%%%%%%%%%%%%%%
\subsection{Inferring Systems Thinking in Problem-Solving}

Within the hierarchical ST framework, \textsc{BEHAVE} can be used to hypothesize candidate ST indicators for further investigation. 
The experts posited, for example, that employing trial-and-error strategies, such as combining \textit{same-start} and \textit{same-end} in a low-sparsity sequence, could indicate low ST levels.
While the log data alone cannot establish whether a pupil is capable of higher-level ST, combining it with dedicated ST measures could test this hypothesis.

A second example was motivated by clustering based on \textit{correctness} and \textit{within-category} rate. 
This analysis yielded four clusters, one of which was characterized both by the lowest average \textit{correctness} score and the lowest \textit{within-category} usage. 
In contrast, the other three clusters exhibited higher average \textit{correctness} and \textit{within-category} rates, together with higher average \textit{tracing} and \textit{cycle} rates. 
Drawing on their domain knowledge, the experts suggested this could reflect difficulties among some pupils in discerning the boundaries of carbon-cycle subsystems, since identifying system boundaries is an important aspect of ST~\cite{monat2015what}. 
In turn, they interpreted this as tentative support for a teaching strategy that focuses on one category at a time.
Although these interaction patterns do not confirm the interpretation in itself, they suggest that focusing on one category at a time may facilitate the emergence of \textit{tracing} and \textit{cycle} strategies associated with higher ST skills.
Within the TC environment, understanding how such ST abilities develop may benefit from cross-task analysis to discern any transfer of acquired knowledge and skills from the terrestrial carbon cycle to the global and ocean carbon cycles.

A further hypothesis rests on the premise that linking reservoirs through processes could be seen as representative of a Synthesis-level ST skill. 
For the participating pupils, a traditional terrestrial food web would be familiar from biology class, whereas the same food chain is less likely to be encountered in an aquatic ecosystem. 
Generalizing knowledge of terrestrial food chains to an aquatic ecosystem could therefore indicate an ST skill at the Implementation level. 
For example,~\cref{fig:two_food} displays data from a pupil who first completes a correct terrestrial food chain, followed shortly thereafter by the corresponding correct food chain in the ocean. 
While the interaction data alone cannot establish which level of ST a pupil relies on, these examples demonstrate how \textsc{BEHAVE} can support the generation of hypotheses about potential log-file indicators of higher-order ST, which can be validated in future assessment-based studies (see~\cref{sec:discussion}).

\begin{figure}
    \centering
    \includegraphics[
    width=\columnwidth,
    alt={
    }
    ]{figs/two_food.pdf}
    \caption{
    Example of a pupil’s interaction sequence who first completes a correct terrestrial food chain, followed by the food chain in the ocean.
    }
    \label{fig:two_food}
\end{figure}




%%%%%%%%%%%%%%%%%%%%%%%%%%%%%%%%%%%%%%%%%%%%%%%%%%%%%%%%%%%%%%%%
% VALIDATION
%%%%%%%%%%%%%%%%%%%%%%%%%%%%%%%%%%%%%%%%%%%%%%%%%%%%%%%%%%%%%%%%
\section{Expert Validation}
\label{sec:expert_validation}

To strengthen the evidential basis of \textsc{BEHAVE}, and building on the case study (\cref{sec:case_study}), we position this section as an \textit{expert validation} grounded in triangulation with prior published results from pupils' interaction with TC~\cite{host2026pupils}, rather than as a new controlled evaluation study.
The domain experts who participated in the exploratory workshops and validation are the interdisciplinary team members introduced in~\cref{sec:design_process}.
They are co-authors of this work, designers of TC, and have previously published an empirical study of the same dataset analyzed here~\cite{host2026pupils}.
Because they were involved both in developing the learning environment and in conducting the prior analysis, the present validation should be understood as a triangulation against established findings rather than an independent evaluation. We return to this distinction when delineating the scope of the validated contribution below.

\noindent
\textbf{Validation Procedure}.
The expert validation was conducted across the exploratory workshops that concluded our design process (\cref{sec:design_process}). 
In these sessions, the experts examined the same $97$-pupil dataset using \textsc{BEHAVE} and were asked, first, whether the system reproduced the strategy structure, error patterns, and task-difficulty effects established in their prior study~\cite{host2026pupils}, and second, what additional observations the system made available beyond that prior analysis.
Sessions combined open-ended exploration with think-aloud commentary.
The experts' judgments, points of agreement with prior findings, and newly emergent observations were captured in session notes and subsequently consolidated by the authors into the validation claims reported below.
Since the previous study analyzed the same interactive TC environment and the same pupils, and crucially, replicated several behavioral patterns that \textsc{BEHAVE} should be able to reveal if it is analytically valid, it provides a meaningful baseline for validation.

The prior study~\cite{host2026pupils} reported that pupils' arrow-drawing behavior emerged as four strategies, namely \textit{Tracing}, \textit{Minimal Sub-Cycle}, \textit{Exhausting Options}, and \textit{Focused Completion}. These strategies varied across arrow-drawing tasks, and consistent error patterns occurred, including frequent incorrect connections between plants and decomposers.
Notably, the most complex task involving $12$ reservoirs (\cref{fig:tcInterface}) was associated with the highest error rate.
Read against this baseline, expert validation of \textsc{BEHAVE} concerns two questions.
First, can the VA system \textit{reproduce core findings} already established by prior statistical and log-based analyses?
Second, can it \textit{go beyond those findings} by making their sequential and temporal structure, and the variation in within-task interactions, visible in ways that were previously inaccessible?
This framing is important for interpreting the contribution of the present work.
\textsc{BEHAVE} is not a replacement for prior analyses, but a visual analytical medium that exposes both known and previously inaccessible patterns at the motif, cohort, and individual levels, while supporting the formulation of more fine-grained hypotheses about systems-thinking-related behavior.

\noindent
\textbf{Reproducing established findings}.
In previous work, the \textit{Exhausting Options} strategy referred to repeatedly drawing arrows from one reservoir to several others, while \textit{Focused Completion} referred to repeatedly drawing arrows from several reservoirs toward one reservoir~\cite{host2026pupils}.
The experts recognized these same strategies in \textsc{BEHAVE}, where they are formalized as the ST-related event dimensions \textit{same-start} and \textit{same-end}, respectively, and appear as clearly separate profiles in the \textsc{Cluster Panel} and \textit{Comparison View}.
Likewise, the experts identified the previously reported \textit{Tracing} and \textit{Minimal Sub-Cycle} strategies as directly represented, appearing as recurring motifs rather than only as overall averages.
Beyond confirming their presence, the experts observed that the cluster information in \textsc{BEHAVE} reveals how these strategies are distributed across different cohorts (e.g.,~\cref{fig:four_clusters}), a level of detail not perceivable in the previous analysis.

\noindent
\textbf{Extending established findings}.
Beyond confirmation, the experts judged that \textsc{BEHAVE} improves the precision of prior observations on task difficulty and problem-solving approaches.
The prior study~\cite{host2026pupils} showed that pupils made more errors on more complex arrow-drawing tasks and suggested that strategies such as Exhausting options might reflect a more mechanical or trial-and-error approach as task demands increase.
With \textsc{BEHAVE}, the experts found this interpretation became far more salient: low-sparsity, rapidly repeated alternations of \textit{tracing} and \textit{mini-cycle} moves could be seen directly in complete sequences, motif matches, and individual trajectories.
What was previously a statistically informed suspicion of trial-and-error could thus be examined by the experts as a concrete, localized behavioral pattern (e.g.,~\cref{fig:trial_error}) with identifiable motif structure and distribution across pupils.

A further aspect concerns domain relevance.
In the prior work~\cite{host2026pupils}, the logs and task outcomes were interpreted in relation to conceptual and reasoning challenges in learning about the carbon cycle, including incorrect links between decomposers and plants and difficulties tracing carbon movement through subsystems.
The experts confirmed that \textsc{BEHAVE} supports the same interpretations, but does so by connecting aggregate correctness to previously inaccessible behaviors.
For instance, the case study showed how pupils may begin with an extended \textit{tracing} move, then shift to \textit{same-start} exploration, or repeated \textit{mini-cycle} between local back-and-forth moves.
These views allowed the experts to relate potential conceptual difficulties and errors not only to outcome results, but to the fine-grained sequence of actions through which those results emerged.
Whereas the earlier study characterized pupils' behavior primarily through task-level statistics, such as initial actions, error frequencies, and average rates, the experts noted that \textsc{BEHAVE} enables them to inspect how strategies are combined within a single interaction history, when they appear, how long they persist, and which motifs recur across pupils.
In this way, the experts confirmed that \textsc{BEHAVE} goes beyond the baseline.
Pupils do not simply employ a single strategy type but often shift between strategies during problem solving (\cref{fig:switch_strategies}), and local structures can be shared across pupils even when complete sequences differ, aspects that the prior statistical analysis highlighted but could not directly visualize.

\noindent
\textbf{Scope of the validated contribution}.
The expert validation supports a measured claim about \textsc{BEHAVE}.
The system should not be seen as ``automatically diagnosing'' ST, nor as a replacement for pedagogical assessment. 
Notably, the expert validation concerns the behavioral patterns that \textsc{BEHAVE} exposes and their correspondence to prior findings. It does not validate the further mapping of these behaviors onto specific ST hierarchy levels, which would require other assessment data (e.g., pre- and post-ST tests or think-aloud protocols) to confirm.
Its validated value therefore lies in enabling experts to actively and collaboratively connect prior results to richer visual evidence of pupils' progression and behavior in TC, thereby making interaction patterns more systematically visible and interpretable.
While the prior study established that strategy patterns exist in the data~\cite{host2026pupils}, the experts confirmed that \textsc{BEHAVE} helps them see where, when, and how those patterns unfold, providing a foundation for hypothesis refinement and for future integration with direct measures of ST.





%%%%%%%%%%%%%%%%%%%%%%%%%%%%%%%%%%%%%%%%%%%%%%%%%%%%%%%%%%%%%%%%
% DISCUSSION
%%%%%%%%%%%%%%%%%%%%%%%%%%%%%%%%%%%%%%%%%%%%%%%%%%%%%%%%%%%%%%%%
\section{Discussion}
\label{sec:discussion}

\textsc{BEHAVE} demonstrates how coupling low-level interaction logs with an ST framework can reveal previously hidden behavioral patterns and support the generation of hypotheses about pupils' reasoning processes. 
Below, we reflect on broader lessons learned from the design and evaluation process, the approach's generalizability and scalability, its current limitations, and promising future directions.

\noindent
\textbf{Lessons learned and pedagogical implications}.
A central lesson is the value of an iterative, co-design process that actively involves domain experts and visualization researchers. 
Early workshops revealed that aggregate statistics and manual inspection of selected logs were limited for direct interpretation regarding ST.
By deriving a set of interpretable behavioral descriptors and presenting them through an ST lens, we were able to move from low-level interactions through ST-related events to recurring patterns, cohort-level tendencies, and candidate interpretations of pupil reasoning.
In this sense, the abstraction layer was not merely a preprocessing step, but a necessary bridge between logged behavior and educationally meaningful analysis.

Moreover, the interdisciplinary approach employed, grounded in data from a learning environment and informed by pedagogical and domain‑specific principles, revealed a need for stronger dialogue between education domain experts and visualization researchers.
The TC environment was intentionally designed to support understanding of complex learning processes.
The co-design process showed that establishing interdisciplinary common ground requires structured support from the outset.
Since educational data are inherently complex, dynamic, influenced by multiple interacting factors, and often challenging to interpret, their analysis demands specialized skills among all involved stakeholders, as well as a shared understanding of methodological and interpretive frameworks.

Finally, the case study demonstrates that ST can serve as an \textit{operational} interpretive lens rather than merely a learning goal, and it suggests several pedagogical implications.
First, since the identified ST-related strategy profiles were associated with different correctness outcomes, the system can help distinguish between productive and less productive forms of systematic exploration, thereby supporting more targeted instructional guidance.
Second, the detection of rapid and repeated \textit{mini-cycle} behavior suggests that some pupils may engage in trial-and-error with limited time for reflection, pointing to opportunities for pedagogical scaffolding that encourages pausing, planning, and explaining transitions between reservoirs. 
Reflecting on these observations during the exploratory sessions, the experts considered such distinctions to be pedagogically actionable. In their view, knowing not only that a pupil struggled, but also \textit{when} and \textit{how} a particular strategy unfolded, could help direct instructional support to where it is most needed, guidance that their prior %aggregate-statistics 
workflow could not provide.
Hence, the pedagogical value of the approach lies in supporting educators and researchers in formulating evidence-based hypotheses about where pupils may need support and which interventions may encourage deeper ST reasoning. 

\noindent
\textbf{Transferability}.
Although \textsc{BEHAVE} was developed for the \textit{Tracing Carbon} environment, the broader analytical approach is not limited to it. 
Our proposed approach is conceptually generalizable beyond ST in the carbon-cycle domain. 
What is transferable is the principle of coupling theory-informed abstractions with multi-level sequence analysis.
In another learning context, the specific descriptors and interpretive constructs would differ.
Still, the same design logic could support the investigation and hypothesis generation across interaction and educational frameworks and theories. 
Thus, the contribution is not only a VA system for a specific analysis, but a workflow for connecting interaction data to domain-relevant constructs.

The ST-framework lens could also be generalized to other complex-system domains in science education, such as climate modeling, ecosystem dynamics, or biochemical networks. 
In principle, any learning environment whose interactions can be mapped to relational events (e.g., connections, cycles, feedback) can adopt the same progression from event to pattern to structure (\textbf{R1}--\textbf{R3}). 
Future adaptations would primarily involve redefining behavioral descriptors and adjusting motif-ranking metrics to domain-specific reasoning constructs.

\noindent
\textbf{Scalability}.
Beyond transferability across domains, an important practical aspect concerns how \textsc{BEHAVE} scales as the number of pupils and length of interaction sequences increases, raising distinct visual and computational considerations.
On the visual side, the \textit{Motif Library} does not grow without bound as the number of pupils increases.
Because the analyzed task has a finite scope, there is a limited number of meaningful ways to solve the carbon cycle, and the discovered motifs capture recurring \textit{shared} behavioral patterns rather than per-pupil sequences, so additional pupils predominantly reinforce existing motifs rather than introduce new ones.
The size of the library is further controllable through the motif-discovery thresholds described in~\cref{subsec:data_abstraction} (e.g., the match-distance threshold $\tau$ and the merge tolerance $\varepsilon$), which directly govern how finely behaviors are distinguished and merged, allowing the analyst to balance granularity against compactness.

By design, the \textsc{Individual Panel} is devoted to a single selected sequence and is therefore unaffected by the number of pupils.
A visual scalability concern at the individual level is instead sequence \textit{length}.
For long interaction sequences, the Marey chart can become dense, which could be addressed in future work through abstraction or summarization techniques, such as semantic zoom or segment-level aggregation.

On the computational side, motif discovery is an offline, one-time preprocessing step that is decoupled from interactive use,  not affecting responsiveness during analysis.
Its dominant cost is the pairwise matching of pooled behavioral windows, which scales with the total number of windows (and thus with both pupil count and sequence length).
This cost is mitigated in our pipeline by activity-pool pruning ($\theta_{\mathrm{pool}}$) and per-pupil deduplication, which reduces the candidate set.
Our approach builds on matrix-profile and consensus-motif methods~\cite{kamgar2019matrix, yeh2017matrix}, whose established acceleration techniques, such as GPU-based variants~\cite{zhu2016matrix}, enable substantially larger datasets.

\noindent
\textbf{Additional Data Sources}.
The integration of richer, multi-modal data, which could substantially deepen the analysis our current log-only setting permits, could strengthen our approach.
Think-aloud protocols are particularly valuable because verbalized reasoning offers one of the closest available windows into cognitive processes that interaction logs can only indirectly reflect, thereby supporting the validation of ST-related hypotheses that the present work is only able to generate.
Eye-tracking could expose the attentional and planning behavior that precedes action, helping to distinguish, for instance, whether the low-sparsity patterns we associate with trial-and-error reflect genuine impulsivity or rapid yet deliberate execution.
Screen recordings could capture hesitations, partial drags, and abandoned actions that were not logged, thereby surfacing micro-level behavior.

These additional data sources were not applicable to this work for reasons beyond the added preprocessing computation.
Our contribution specifically targets the interaction logs that TC already produces at scale in grade $7$--$9$ classrooms.  Capturing gaze, speech, or video would alter the data-collection apparatus and study design, and also require stricter consent requirements, especially in cases where participants are younger than 15 years of age.
Moreover, the principal challenge is not data preprocessing, but the temporal alignment and merging of modalities with differing sampling rates and semantics, as well as the interpretation of their interrelationships, which constitutes a research problem in its own right.
Our architecture is nonetheless well-positioned for such an extension. Additional modalities could be encoded as further behavioral descriptors that feed into the same pipeline (\textbf{R1}--\textbf{R3}), with the sliding-window abstraction serving as a natural temporal substrate for aligning gaze, speech, and interaction events.
We regard this multi-modal extension as a primary avenue for future work, particularly in conjunction with pre- and post-ST assessments discussed above.

\noindent
\textbf{Limitations and Future Work}.
While \textsc{BEHAVE} demonstrates several valuable insights, its current implementation has limitations.
Specifically, our evaluation and validation were restricted to a single complex task within one digital environment. Although it drew on data from $97$ pupils, broader generalization would benefit from extended studies.
While the case study and expert validation yielded rich qualitative interpretations, larger-scale, multi-task, or multi-site studies could be beneficial for further evaluating the robustness of our approach.
We note that this qualitative and expert-focused form of evaluation is a deliberate methodological choice rather than a substitute for a controlled study.
Evaluation paradigms based on task-completion time, accuracy, or benchmark comparisons assume well-defined tasks with an established ground truth. In contrast, \textsc{BEHAVE} targets open-ended, exploratory hypothesis generation about pupils' ST reasoning, an inherently unobservable construct for which no such ground truth is available.
Consistent with established evaluation approaches for visual analytics~\cite{lam2012empirical, carpendale2008evaluating} and design-study methodology~\cite{sedlmair2012design}, an insight-oriented expert evaluation is appropriate for assessing a tool of this nature, and we report our findings accordingly.
Supporting this framing, \textsc{BEHAVE} is designed to help analysts hypothesize and explore links between motifs and ST hierarchy levels rather than to confirm them. Consequently, it does not provide ground-truth validation of pupils' ST reasoning.
Future integration of pre- and post-ST assessments, together with other data, could enable quantitative mapping of behavioral indicators to cognitive constructs, thus extending rather than replacing exploratory analyses supported by \textsc{BEHAVE}.





%%%%%%%%%%%%%%%%%%%%%%%%%%%%%%%%%%%%%%%%%%%%%%%%%%%%%%%%%%%%%%%%
% CONCLUSION
%%%%%%%%%%%%%%%%%%%%%%%%%%%%%%%%%%%%%%%%%%%%%%%%%%%%%%%%%%%%%%%%
\section{Conclusion}

We have presented a visual analytics approach that integrates multidimensional time-series representation, motif discovery, and multi-level visual representations to support the exploration of pupils' problem-solving behaviors in a digital learning environment through an ST lens.
The approach is implemented in \textsc{BEHAVE} and was developed through an iterative co-design process with domain experts using interaction log data from the \textit{Tracing Carbon} environment. 
Throughout this process, collaborative workshops shaped various key design decisions, including the dual view of sequential and durational interaction, the category-grouped representation, and the shift from precomputed communities to interactive clustering.
Our results demonstrate the usefulness and validity of the approach for uncovering previously unseen behavioral patterns and generating hypotheses about the relationship between observable interactions and ST reasoning.
We further hope that both \textsc{BEHAVE} and the documented co-design process can inform future efforts to connect interaction data with educational reasoning frameworks.





%%%%%%%%%%%%%%%%%%%%%%%%%%%%%%%%%%%%%%%%%%%%%%%%%%%%%%%%%%%%%%%%
% Supplemental
%%%%%%%%%%%%%%%%%%%%%%%%%%%%%%%%%%%%%%%%%%%%%%%%%%%%%%%%%%%%%%%%
\section*{Supplemental Materials}
\label{sec:supplemental_materials}

All supplemental materials are available on OSF at \url{https://osf.io/2skm8/}, released under a CC BY 4.0 license.





\section*{Acknowledgments}
The work was supported by the Swedish Research Council (grant 2020-05000) and, in part, by the Swedish Research Council (grant 2020-05147), the Marcus and Amalia Wallenberg Foundation (2023.0130), and the Knut and Alice Wallenberg Foundation (2019.0024).
Generative AI tools were used for language editing, including improvements to grammar, phrasing, and clarity. 
The authors reviewed and revised all suggestions and take full responsibility for the final text. 


\bibliographystyle{IEEEtran} 
\bibliography{template}


\newpage

 
\vspace{11pt}


\vspace{11pt}


\vfill

\end{document}


